% templates/pdf-header.tex
% Styling cho PDF export — Cười Vỡ Bụng
% Dùng với LuaLaTeX + pandoc --include-in-header

% ============================================================
% EMOJI FONT FALLBACK
% Đăng ký Noto Color Emoji làm fallback cho ký tự emoji/icon.
% LuaLaTeX + luaotfload tự động chuyển sang font này khi gặp
% ký tự mà Noto Sans không có.
% ============================================================
\directlua{
  luaotfload.add_fallback("emojifallback", {
    "NotoColorEmoji:mode=harf;"
  })
}
\setmainfont{Noto Sans}[
  RawFeature={fallback=emojifallback},
  BoldFont={Noto Sans Bold},
  ItalicFont={Noto Sans Italic},
  BoldItalicFont={Noto Sans Bold Italic}
]

% ============================================================
% MÀU SẮC
% ============================================================
\definecolor{chaptercolor}{HTML}{1A5276}
\definecolor{storycolor}{HTML}{2E4053}
\definecolor{quoteborder}{HTML}{2980B9}
\definecolor{quotebg}{HTML}{EBF5FB}
\definecolor{rulegray}{HTML}{AEB6BF}

% ============================================================
% BLOCKQUOTE STYLING
% Pandoc render markdown > thành \begin{quote}...\end{quote}.
% Thay vì indent nhạt mặc định, dùng tcolorbox: viền trái xanh,
% nền nhạt — trông giống callout trong EPUB/web.
% ============================================================
\usepackage{tcolorbox}
\tcbuselibrary{breakable,skins}

\renewenvironment{quote}{%
  \begin{tcolorbox}[
    enhanced,
    breakable,
    colback=quotebg,
    colframe=quoteborder,
    leftrule=3.5pt,
    rightrule=0pt,
    toprule=0pt,
    bottomrule=0pt,
    boxsep=2pt,
    left=10pt,
    right=10pt,
    top=6pt,
    bottom=6pt,
    arc=0pt,
    outer arc=0pt,
  ]%
}{%
  \end{tcolorbox}%
}

% ============================================================
% HEADING STYLING
% section = H1 (chương), subsection = H2, subsubsection = H3 (truyện)
% ============================================================
\usepackage{titlesec}

\titleformat{\section}
  {\LARGE\bfseries\color{chaptercolor}}
  {}{0em}{}[\vspace{0.3em}{\color{rulegray}\titlerule[1pt]}]
\titlespacing*{\section}{0pt}{2em}{1em}

\titleformat{\subsection}
  {\Large\bfseries\color{chaptercolor}}
  {}{0em}{}
\titlespacing*{\subsection}{0pt}{1.5em}{0.8em}

\titleformat{\subsubsection}
  {\large\bfseries\color{storycolor}}
  {}{0em}{}
\titlespacing*{\subsubsection}{0pt}{1.2em}{0.5em}

% ============================================================
% HORIZONTAL RULE — tô màu xám nhạt thay vì đen
% ============================================================
\let\OldRule\rule
\renewcommand{\rule}[2]{\textcolor{rulegray}{\OldRule{#1}{#2}}}

% ============================================================
% HEADER / FOOTER
% ============================================================
\usepackage{fancyhdr}
\pagestyle{fancy}
\fancyhf{}
\fancyhead[C]{\small\color{rulegray}\textit{Cười Vỡ Bụng — Tuyển Tập Truyện Cười Thế Giới Cho Người Việt}}
\fancyfoot[C]{\small\color{storycolor}\thepage}
\renewcommand{\headrulewidth}{0.4pt}
\renewcommand{\headrule}{\hbox to\headwidth{\color{rulegray}\leaders\hrule height \headrulewidth\hfill}}

% Trang đầu mỗi chương dùng plain — chỉ số trang, không header
\fancypagestyle{plain}{%
  \fancyhf{}%
  \fancyfoot[C]{\small\color{storycolor}\thepage}%
  \renewcommand{\headrulewidth}{0pt}%
}

% ============================================================
% KHOẢNG CÁCH & LINE SPACING
% ============================================================
\setlength{\parskip}{0.5em}
\setlength{\parindent}{0em}

\usepackage{setspace}
\setstretch{1.25}
